En ce qui concerne la friction lors de la descente selon l'ouverture de la valve,
deux d'approximations  sont faites (linéaire et quadratique). La méthode des
moindres carrés est utilisée pour déterminer la courbe passant le plus proche
des points selon l'ordre choisi (linéaire étant un ordre 1). Seul le
nombre d'inconnus/coefficients $a$ diffère dans la démarche pour les deux
approximations. La plus complexe est donnée en exemple ici:

{\footnotesize
\begin{align}
	Y = \begin{bmatrix} 0.87\\0.78\\0.71\\0.61\\0.62
	\\0.51\\0.51\\0.49\\0.46\\0.48\\0.46 \end{bmatrix}\hspace{.5cm}
		M &=
	\begin{bmatrix}
	1 &0   &0  \\
	1 &10  &10^2  \\
	1 &20  &20^2  \\
	1 &30  &30^2  \\
	1 &40  &40^2  \\
	1 &50  &50^2  \\
	1 &60  &60^2  \\
	1 &70  &70^2  \\
	1 &80  &80^2  \\
	1 &90  &90^2  \\
	1 &100 &100^2 \\
	\end{bmatrix}\hspace{.5cm}
	V =
	\begin{bmatrix} a_0\\a_1\\a_2 \end{bmatrix} \\
	M^T MV = M^T Y \Rightarrow
	V &= \left( M^T M \right)^{-1} M^T Y \Rightarrow
	V =
	\begin{bmatrix} 0.8668\\-0.0092\\0.0001 \end{bmatrix} \\
		F(x) &= 0.8668 -0.0092x + 0.0001 x^2
		\label{eq:valve-quad}
\end{align}}

Pour chacune des approximations, il est possible de connaitre l'erreur RMS en
faisant la racine de la somme des écarts au carrés (écart entre la valeur
expérimentale et celle prédite par la fonction développée):

\begin{equation}
	\textrm{RMS} = \sqrt{\frac{1}{N}\sum_{i=1}^{N}\Delta_y^2}
\end{equation}
\begin{equation}
	\textrm{RMS}_\textrm{linéaire} = 0.0494\hspace{.5cm}
	\textrm{RMS}_\textrm{quadratique} = 0.0180
\end{equation}

Puisque le RMS est déjà faible pour une approximation quadratique (ordre 2) et
que seules 11 données sont disponibles de 0 à 100\%, il n’est pas pertinent
d’utiliser des polynômes d’ordre supérieur (3, 4, \ldots), qui pourraient
surparamétriser l’approximation, amplifier les erreurs expérimentales ou générer
des oscillations non-représentatives de la réalité. Les deux courbes sont
présentées à la figure \ref{fig:02-valve}, montrant que la plage de coefficients
de friction possible pour le système varie approximativement entre 0,47 et 0,87.

